\section{Diagnosis: what \vfive{} does wrong, and what it does not}
\label{sec:diagnosis}

The v0.5 study closed with an unusually candid non-claim: removing a failure mode was worth almost
nothing in win rate, and that was the result. This section asks the obvious follow-up question ---
where, then, is the remaining value --- and answers it with four measurements, two of which are
negative and are reported as findings rather than as caveats. The negative pair is load-bearing:
between them they close off the two largest items on the design register this study inherited, and
they are the reason \vsix{} spends its effort on the measurement apparatus rather than on the
policy's feature set.

\subsection{The instrument}
\label{sec:diag-instrument}

Every claim below is produced by \texttt{fish v6probe}, a diagnostics subcommand added for this
study, and is reproducible from a clean build by the commands in Appendix~\ref{app:repro}. Three
definitions fix what is being counted.

\begin{definition}[Live candidate]
An ask $(c,t)$ is \emph{live} for an actor if it is legal and the actor cannot prove from its own
hand and the public record that $t$ does not hold $c$. This is the M1 predicate \vfive{} already
gates on; it involves no posterior and no policy prior, so a candidate set restricted to it cannot
have been narrowed by a modelling error.
\end{definition}

\begin{definition}[Exact tie at the top]
A decision has an \emph{exact tie at the top} when two or more live candidates receive scores that
are equal to within $10^{-9}$ under the policy's own pre-search score, that is
$\lambda\, w^{\!\top}\! f + \nu\, Q$ where $Q$ is the one-ply expectimax over the value function.
The gap and the spread are reported alongside so that a tie can be distinguished from a flat
objective.
\end{definition}

\begin{definition}[Belief-symmetric tie]
A tie is \emph{belief-symmetric} when the tied candidates additionally receive equal probability
under the \emph{exact} posterior of \S\ref{sec:belief}, not merely under the Sinkhorn fit the
deployed policy runs. A belief-symmetric tie is a statement about the information set, not about
the approximation.
\end{definition}

\subsection{More than half of all ask decisions end in a tie}
\label{sec:diag-ties}

Over \vsixTieNFive{} tied decisions in \vfive{} mirror play, \vsixTieShareFive\% of contested
decisions --- decisions with at least two live candidates, which is \vsixPathFiveContested\% of all
of them --- carry an exact tie at the top. The objective is not flat: the mean spread across the
candidate set is \vsixTieSpreadFive{} while the mean gap between the leaders is
\vsixTieGapFive{}. The dimension is blind, not the surface.

The tie is structural, and it is structural in a specific place. \vsixTieSameSuitFive\% of ties are
two \emph{cards of one half-suit at one target}; only \vsixTieSameCardFive\% are one card offered to
two different targets. Three causes stack. Eighteen of the twenty ask features are computed per
(half-suit, target) and the three per-card features are all functions of the same marginal; the
Sinkhorn filter makes exchangeable cards numerically identical by construction; and
\texttt{askExpectedValue} discards its \texttt{target} argument outright, so the expectimax half of
the score is constant across targets. The last of these was the defect the v0.5 design register
headlined. It is aimed at \vsixTieSameCardFive\% of the tie mass.

\subsection{Negative result: the ties are exchangeable, and therefore irreducible}
\label{sec:diag-irreducible}

The natural reading of \S\ref{sec:diag-ties} is that a subroutine \vfive{} already runs is throwing
away a large amount of value, and the register this study inherited priced that value at over three
half-suits per game using a hindsight oracle. That reading is wrong, and the measurement that
refutes it is simple: re-score the tied candidates under the exact count law of
\S\ref{sec:belief-exactworse} and ask how many ties it separates.

\begin{center}
\begin{tabular}{lr}
\toprule
tie-break rule & realised hit rate \\
\midrule
enumeration order (whichever candidate was emitted first) & \vsixTieHitArrayFive\% \\
argmax of the deployed Sinkhorn marginal, within the tie & \vsixTieHitFastFive\% \\
the shipped policy's actual pick (its chain/threat pass) & \vsixTieHitShipFive\% \\
argmax of the \emph{exact} posterior, within the tie & \vsixTieHitExactFive\% \\
\midrule
hindsight best within the tie & \vsixTieHitBestFive\% \\
\bottomrule
\end{tabular}
\end{center}

The exact posterior separates \vsixTieExactSepPctFive\% of the ties, and every rule realises the
same hit rate to three significant figures. Two cards of one half-suit at one target, tied under the
policy's score, are tied under the exact posterior as well. The \vsixTieHitBestFive\% column prices
clairvoyance, not headroom.

\paragraph{The scope of that result, and the experiment that fixes it.} What the table measures is
equality of the \emph{marginal} probability that the target holds the card. It follows that no rule
ranking the tied candidates by that marginal --- the deployed one, the exact one, or any other ---
can beat enumeration order, and none does. An earlier draft of this section concluded from that
that the ties are irreducible. They are not, and the four rows below are the correction. Each plays
the same policy against \vfive{}, changing only how the tie group is resolved.

\begin{center}
\begin{tabular}{lr}
\toprule
how the tie group is resolved & win rate against \vfive{} \\
\midrule
enumeration order (the shipped rule) & $50.00$ by construction \\
uniformly at random, seeded from the public event stream & \vsixSearchTieRandom\% [\vsixSearchTieRandomLo, \vsixSearchTieRandomHi] \\
by a one-deal rollout & \vsixSearchTieOnlyOne\% [\vsixSearchTieOnlyOneLo, \vsixSearchTieOnlyOneHi] \\
\textbf{by a twelve-deal rollout ensemble} & \textbf{\vsixSearchTieOnlyTwelve\%} [\vsixSearchTieOnlyTwelveLo, \vsixSearchTieOnlyTwelveHi] \\
\bottomrule
\end{tabular}
\end{center}

Randomising the choice is worth \emph{exactly} nothing --- $50.00\%$, mean half-suits $4.500$ to
$4.500$ --- so this is not a story about a deterministic policy being predictable. One sampled deal
is worth nothing either. Twelve are worth several points at the bank in the table, and pooled over
every bank the configuration was measured on, \vsixSearchPooled\% over \vsixSearchPooledN{} games
with no bank below \vsixSearchPooledMin\%.

The separating information is therefore real, and it is not the mere fact of varying a
deterministic choice --- the random row rules that out exactly. It is in the \emph{joint}
posterior. Two cards can be equally likely to sit with the target while being differently correlated
with the rest of the deal --- differently placed in the allocations that survive the certificates,
differently useful to the half-suits the opponents are collecting --- and every marginal, exact or
not, integrates that structure away. An ensemble of whole sampled deals does not.

\paragraph{What is not resolved.} Whether the search's advantage is \emph{confined} to the tie
group is a separate question and this study does not answer it. Restricting the search to the tie
group gives \vsixTieSearchBankList\% over \vsixTieSearchBanks{} banks (mean
\vsixTieSearchMean\%); guarding the tie group
instead gives \vsixSearchTieGuardedWin\%; the unrestricted search gives \vsixSearchPooled\%. At
720 games a cell those intervals all overlap. What is established is the pair of negative controls
and the pooled effect, not the attribution. \S\ref{sec:search} says so again where it matters.

Two corrections to the record follow, and are carried into \S\ref{sec:corrections}.

\begin{enumerate}
\item \textbf{Enumeration order does not decide half of all asks.} \vfive{}'s top-$K$ chain and
  threat pass re-scores the leaders and moves the pick at \vsixTieShipMovedFive\% of ties, so array
  order settles roughly a fifth of contested decisions rather than more than half --- and, per the
  table, settles them exactly as well as any alternative.
\item \textbf{The channel is the card dimension, not the target dimension.} The emphasis the v0.5
  register placed on \texttt{(void)target;} is aimed at a twentieth of the tie mass.
\end{enumerate}

The design consequence is stated here and used twice later. Because the tie is a property of the
information set, a rollout ensemble drawn from that information set is symmetric over it too: the
determinizations carry no signal about which of two exchangeable cards to name, only variance. The
tie-break question is therefore not a search question, and \S\ref{sec:search} does not treat it as
one.

\subsection{Negative result: the \vfive{} fit was sampling noise}
\label{sec:diag-fitnoise}

The second negative result concerns the interpretation of \vfive{}'s refit rather than its shipped
configuration. Read from the fit trace, the incumbent objective moves over forty generations by
less than one generation's standard deviation, with an ordinary-least-squares slope whose $t$
statistic is below two, and the gap between the best candidate and the incumbent tracks the expected
maximum of a population of independent draws of exactly that noise.

The cause is visible in the sampler rather than in the objective. One scalar step size was applied
to every coordinate of a vector whose coordinates have ranges spanning two and a half orders of
magnitude, so the search was bang-bang on the bounded probability knobs --- the worst of them
clipped at \vsixFitAClipGenZero\% of generation-zero draws under the old setting --- and effectively
frozen on the unbounded weights. The per-coordinate step-size flag the interface advertised was
parsed and discarded. \S\ref{sec:fitting} is the repair, and \S\ref{sec:results} is what a working
optimiser then finds.

Taken with the independent observation that running the v0.5 \emph{mechanisms} on v0.4's frozen
parameter vector reproduces essentially the whole v0.5 result, the conclusion is that the
v0.4$\rightarrow$v0.5 transition was mechanical and that its parametric machinery never moved. That
is not a small correction: it means the shipped vector this study starts from is v0.4's, and v0.4's
was produced by a search carrying a separate aliasing defect of its own.

\subsection{What is actually still leaking}
\label{sec:diag-channels}

With the tie channel removed from the register and the parametric channel re-opened, the residual
loss channels are smaller than the inherited register suggested, and they are ranked here with the
paired figure each was measured at.

\begin{itemize}
\item \textbf{Asks into a half-suit the actor's own team already owns outright.}
  \vsixPathFiveOwnlocked\% of all asks --- guaranteed misses that donate the turn. M1 cannot see
  them: the enumerator only offers opponents as targets, and no single seat can \emph{prove} a
  teammate holds a card. Pricing the posterior mass the team already carries is the graded
  substitute, and \S\ref{sec:mechanisms} reports that it is worth nothing at the resolution
  \S\ref{sec:protocol} establishes as necessary.
\item \textbf{Declaration errors}, \vsixPathFiveDeclwrong\% of declarations in mirror play. The
  pre-gates that produce them are, as \S\ref{sec:engine} reports, exact: \vsixGateFalseNeg{} false
  negatives in \vsixGateRejected{} rejections. The error is in the allocation, not in the gate.
\item \textbf{The move class M1 deletes.} A provably-dead ask at a chosen opponent is how a human
  blackballs, donates the turn deliberately, or pays for a signal. \S\ref{sec:search} shows that
  \vfive{}'s score cannot price it even when it is offered, and that unbanning it wholesale trades
  a higher win rate for a policy that kills \vsixDeadPathLimit\% of its games at the action limit.
\end{itemize}

None of these is large. That is the diagnosis: \vfive{} sits on a plateau of its own policy class,
and the leverage available to \vsix{} is in the apparatus that measures and fits it, not in another
feature.
